
\documentclass{article}
\usepackage{colortbl}
\usepackage{makecell}
\usepackage{multirow}
\usepackage{supertabular}

\begin{document}

\newcounter{utterance}

\centering \large Interaction Transcript for game `cladder', experiment `full\_v1.5\_default', episode 8085 with selfrefine{-}llama{-}3.1{-}8b{-}instruct.
\vspace{24pt}

{ \footnotesize  \setcounter{utterance}{1}
\setlength{\tabcolsep}{0pt}
\begin{supertabular}{c@{$\;$}|p{.15\linewidth}@{}p{.15\linewidth}p{.15\linewidth}p{.15\linewidth}p{.15\linewidth}p{.15\linewidth}}
    \# & \multicolumn{2}{c}{Player} && \multicolumn{2}{c}{Game Master} \\
    \hline

    \theutterance \stepcounter{utterance}  
    & & & \multicolumn{4}{p{0.6\linewidth}}{
        \cellcolor[rgb]{0.9,0.9,0.9}{
            \makecell[{{p{\linewidth}}}]{
                \texttt{\tiny{[P1$\langle$GM]}}
                \texttt{You are an expert in causal inference. The following question is not a typical commonsense query, but rather a meticulously designed question created by a professor specializing in causal inference, intended to assess the students' mastery of the course content.} \\
\texttt{Imagine a self{-}contained, hypothetical world with only the following conditions, and without any unmentioned factors or causal relationships: Kwox has a direct effect on kwoz. Swoq has a direct effect on kwoz. For those who are not kwox, the probability of kwoz is 35\%. For those who are kwox, the probability of kwoz is 65\%. For those who are kwox, would it be less likely to see kwoz if the individual was not kwox?} \\
\texttt{Start your answer with 'Yes' or 'No', followed by additional reasoning or evidence to support your explanation.} \\
            }
        }
    }
     \\ \\

    \theutterance \stepcounter{utterance}  
    & \multicolumn{4}{p{0.6\linewidth}}{
        \cellcolor[rgb]{0.561,0.737,0.561}{
            \makecell[{{p{\linewidth}}}]{
                \texttt{\tiny{[P1$\rangle$GM]}}
                \texttt{No.} \\
\\ 
\texttt{To answer this question, we need to consider the causal relationships and conditional probabilities given. We are asked to compare the probability of kwoz for individuals who are kwox and those who are not kwox, but under the condition that the individual is not kwox.} \\
\\ 
\texttt{Let's denote the probability of kwoz given that the individual is not kwox as P(kwoz|not kwox). We are given that for those who are not kwox, the probability of kwoz is 35\%, or P(kwoz|not kwox) = 0.35.} \\
\\ 
\texttt{We are also given that for those who are kwox, the probability of kwoz is 65\%, or P(kwoz|kwox) = 0.65.} \\
\\ 
\texttt{However, we are asked to compare the probability of kwoz for those who are kwox but not kwox. This is a counterfactual scenario, and we need to consider the causal relationships between kwox and kwoz.} \\
\\ 
\texttt{Since kwox has a direct effect on kwoz, we can use the causal graph to reason about the counterfactual scenario. The causal graph implies that the probability of kwoz is influenced by the presence of kwox.} \\
\\ 
\texttt{To compare the probability of kwoz for those who are kwox but not kwox, we need to consider the potential outcomes under the counterfactual scenario. If an individual is kwox, the probability of kwoz is 65\%. If an individual is not kwox, the probability of kwoz is 35\%.} \\
\\ 
\texttt{However, the question asks us to compare the probability of kwoz for those who are kwox but not kwox. This is a subtle distinction, as it implies that we are considering the probability of kwoz for individuals who are kwox, but under the condition that they are not kwox.} \\
\\ 
\texttt{In this case, we can use the law of total probability to reason about the counterfactual scenario. Let's denote the probability of kwox as P(kwox). We can then calculate the probability of kwoz for those who are kwox but not kwox as follows:} \\
\\ 
\texttt{P(kwoz|kwox, not kwox) = P(kwoz|not kwox) * P(kwox|not kwox) / P(not kwox) + P(kwoz|kwox) * P(kwox|kwox) / P(kwox)} \\
\\ 
\texttt{However, we are given that P(kwoz|not kwox) = 0.35 and P(kwox|kwox) = 1. We also know that P(kwox) = 1 {-} P(not kwox). } \\
\\ 
\texttt{Since P(kwoz|kwox)} \\
            }
        }
    }
    & & \\ \\

    \theutterance \stepcounter{utterance}  
    & & & \multicolumn{2}{p{0.3\linewidth}}{
        \cellcolor[rgb]{0.9,0.9,0.9}{
            \makecell[{{p{\linewidth}}}]{
                \texttt{\tiny{[GM$|$GM]}}
                \texttt{No} \\
            }
        }
    }
    & & \\ \\

    \theutterance \stepcounter{utterance}  
    & & & \multicolumn{2}{p{0.3\linewidth}}{
        \cellcolor[rgb]{0.9,0.9,0.9}{
            \makecell[{{p{\linewidth}}}]{
                \texttt{\tiny{[GM$|$GM]}}
                \texttt{yes} \\
            }
        }
    }
    & & \\ \\

    \theutterance \stepcounter{utterance}  
    & & & \multicolumn{2}{p{0.3\linewidth}}{
        \cellcolor[rgb]{0.9,0.9,0.9}{
            \makecell[{{p{\linewidth}}}]{
                \texttt{\tiny{[GM$|$GM]}}
                \texttt{game\_result = LOSE} \\
            }
        }
    }
    & & \\ \\

\end{supertabular}
}

\end{document}
